\section{AHSME 1982}
        \begin{enumerate}
        	\item (\textit{AHSME 1982}) When the polynomial $x^3-2$ is divided by the polynomial $x^2-2$, the remainder is 


 $\textbf{(A)} \ 2 \qquad  \textbf{(B)} \ -2 \qquad  \textbf{(C)} \ -2x-2 \qquad  \textbf{(D)} \ 2x+2 \qquad  \textbf{(E)} \ 2x-2$


	\item (\textit{AHSME 1982}) If a number eight times as large as $x$ is increased by two, then one fourth of the result equals 


 $\textbf{(A)} \ 2x + \frac{1}{2} \qquad  \textbf{(B)} \ x + \frac{1}{2} \qquad  \textbf{(C)} \ 2x+2 \qquad  \textbf{(D)}\ 2x+4 \qquad \textbf{(E)}\ 2x+16$


	\item (\textit{AHSME 1982}) Evaluate $(x^x)^{(x^x)}$ at $x = 2$. 


 $\textbf{(A)} \ 16 \qquad  \textbf{(B)} \ 64 \qquad  \textbf{(C)} \ 256 \qquad  \textbf{(D)} \ 1024 \qquad  \textbf{(E)} \ 65,536$


	\item (\textit{AHSME 1982}) The perimeter of a semicircular region, measured in centimeters, is numerically equal to its area, 
measured in square centimeters. The radius of the semicircle, measured in centimeters, is 


 $\textbf{(A)} \ \pi \qquad  \textbf{(B)} \ \frac{2}{\pi} \qquad  \textbf{(C)} \ 1 \qquad  \textbf{(D)} \ \frac{1}{2}\qquad \textbf{(E)} \ \frac{4}{\pi}+2$


	\item (\textit{AHSME 1982}) Two positive numbers $x$ and $y$ are in the ratio $a: b$ where $0 < a < b$. If $x+y = c$, then the smaller of $x$ and $y$ is 


 $\textbf{(A)} \ \frac{ac}{b} \qquad  \textbf{(B)} \ \frac{bc-ac}{b} \qquad  \textbf{(C)} \ \frac{ac}{a+b} \qquad  \textbf{(D)}\ \frac{bc}{a+b}\qquad \textbf{(E)}\ \frac{ac}{b-a}$


	\item (\textit{AHSME 1982}) The sum of all but one of the interior angles of a convex polygon equals $2570^\circ$. The remaining angle is 


 $\textbf{(A)} \ 90^\circ \qquad  \textbf{(B)} \ 105^\circ \qquad  \textbf{(C)} \ 120^\circ \qquad  \textbf{(D)}\ 130^\circ\qquad \textbf{(E)}\ 144^\circ$


	\item (\textit{AHSME 1982}) If the operation $x \star y$ is defined by $x \star y = (x+1)(y+1) - 1$, then which one of the following is FALSE? 


 $\textbf{(A)} \ x \star y = y\star x  \text{ for all real } x,y. \\ \textbf{(B)} \ x \star (y + z) = ( x \star y ) + (x \star z)  \text{ for all real } x,y, \text{ and } z.\\ \textbf{(C)} \ (x-1) \star (x+1) = (x \star x) - 1 \text{ for all real } x. \\ \textbf{(D)} \ x \star 0 = x \text{ for all real } x. \\ \textbf{(E)} \ x \star (y \star z) = (x \star y) \star z \text{ for all real } x,y, \text{ and } z.$


	\item (\textit{AHSME 1982}) By definition, $r! = r(r - 1) \cdots 1$ and $\binom{j}{k} = \frac {j!}{k!(j - k)!}$, where $r,j,k$ are positive integers and $k < j$. 
If $\binom{n}{1}, \binom{n}{2}, \binom{n}{3}$ form an arithmetic progression with $n > 3$, then $n$ equals 


 $\textbf{(A)}\ 5\qquad  \textbf{(B)}\ 7\qquad  \textbf{(C)}\ 9\qquad  \textbf{(D)}\ 11\qquad  \textbf{(E)}\ 12$


	\item (\textit{AHSME 1982}) A vertical line divides the triangle with vertices $(0,0), (1,1)$, and $(9,1)$ in the $xy\text{-plane}$ into two regions of equal area. 
The equation of the line is $x=$
$\textbf {(A)}\ 2.5 \qquad  \textbf {(B)}\ 3.0 \qquad  \textbf {(C)}\ 3.5 \qquad  \textbf {(D)}\ 4.0\qquad  \textbf {(E)}\ 4.5$


	\item (\textit{AHSME 1982}) In the adjoining diagram, $BO$ bisects $\angle CBA$, $CO$ bisects $\angle ACB$, and $MN$ is parallel to $BC$. 
If $AB=12, BC=24$, and $AC=18$, then the perimeter of $\triangle AMN$ is


 
                    \begin{figure}[H]
                        \centering
                        \includegraphics[width=0.5\textwidth]{Problem/AHSME_I_1982_10_Problem_Figure_1.png}
                    \end{figure}
                    
$\textbf {(A)}\ 30 \qquad  \textbf {(B)}\ 33 \qquad  \textbf {(C)}\ 36 \qquad  \textbf {(D)}\ 39 \qquad  \textbf {(E)}\ 42$


	\item (\textit{AHSME 1982}) How many integers with four different digits are there between $1,000$ and $9,999$ such that the absolute value of 
the difference between the first digit and the last digit is $2$? 


 $\textbf {(A)}\ 672 \qquad  \textbf {(B)}\ 784 \qquad  \textbf {(C)}\ 840 \qquad  \textbf {(D)}\ 896 \qquad  \textbf {(E)}\ 1008$


	\item (\textit{AHSME 1982}) Let $f(x) = ax^7+bx^3+cx-5$, where $a,b$ and $c$ are constants. If $f(-7) = 7$, then $f(7)$ equals 


 $\textbf {(A)}\ -17 \qquad  \textbf {(B)}\ -7 \qquad \textbf {(C)}\ 14 \qquad  \textbf {(D)}\ 21\qquad  \textbf {(E)}\ \text{not uniquely determined}$


	\item (\textit{AHSME 1982}) If $a>1, b>1$, and $p=\frac{\log_b(\log_ba)}{\log_ba}$, then $a^p$ equals 


 $\textbf {(A)}\ 1 \qquad  \textbf {(B)}\ b \qquad  \textbf {(C)}\ \log_ab \qquad  \textbf {(D)}\ \log_ba \qquad  \textbf {(E)}\ a^{\log_ba}$


	\item (\textit{AHSME 1982}) In the adjoining figure, points $B$ and $C$ lie on line segment $AD$, and $AB, BC$, and $CD$ are diameters of circle 
$O, N$, and $P$, respectively. Circles $O, N$, and $P$ all have radius $15$ and the line $AG$ is tangent to circle $P$ at $G$. 
If $AG$ intersects circle $N$ at points $E$ and $F$, then chord $EF$ has length


 
                    \begin{figure}[H]
                        \centering
                        \includegraphics[width=0.5\textwidth]{Problem/AHSME_I_1982_14_Problem_Figure_1.png}
                    \end{figure}
                    
$\textbf {(A)}\ 20 \qquad  \textbf {(B)}\ 15\sqrt{2} \qquad  \textbf {(C)}\ 24 \qquad  \textbf {(D)}\ 25 \qquad  \textbf {(E)}\ \text{none of these}$


	\item (\textit{AHSME 1982}) Let $[z]$ denote the greatest integer not exceeding $z$. Let $x$ and $y$ satisfy the simultaneous equations 


 
 \begin{align*} y&=2[x]+3 \\ y&=3[x-2]+5. \end{align*}



 If $x$ is not an integer, then $x+y$ is 


 $\textbf{(A) } \ \text{ an integer} \qquad  \textbf{(B) } \ \text{ between 4 and 5} \qquad  \textbf{(C) } \ \text{ between  -4 and 4}\qquad\\ \textbf{(D) } \ \text{ between 15 and 16}\qquad \textbf{(E) } \  16.5$


	\item (\textit{AHSME 1982}) A wooden cube has edges of length $3$ meters. Square holes, of side one meter, centered in each face are cut through to the opposite face. 
The edges of the holes are parallel to the edges of the cube. The entire surface area including the inside, in square meters, is 


 $\textbf {(A)} \ 54 \qquad  \textbf {(B)} \ 72 \qquad  \textbf {(C)} \ 76 \qquad  \textbf {(D)} \ 84\qquad  \textbf {(E)} \ 86$


	\item (\textit{AHSME 1982}) How many real numbers $x$ satisfy the equation $3^{2x+2}-3^{x+3}-3^{x}+3=0$? 


 $\textbf {(A)}\ 0 \qquad  \textbf {(B)}\ 1 \qquad  \textbf {(C)}\ 2 \qquad  \textbf {(D)}\ 3 \qquad  \textbf {(E)}\ 4$


	\item (\textit{AHSME 1982}) In the adjoining figure of a rectangular solid, $\angle DHG=45^\circ$ and $\angle FHB=60^\circ$. Find the cosine of $\angle BHD$.


 
                    \begin{figure}[H]
                        \centering
                        \includegraphics[width=0.5\textwidth]{Problem/AHSME_I_1982_18_Problem_Figure_1.png}
                    \end{figure}
                    
$\textbf{(A)}\ \frac{\sqrt{3}}{6} \qquad  \textbf{(B)}\ \frac{\sqrt{2}}{6} \qquad  \textbf{(C)}\ \frac{\sqrt{6}}{3} \qquad  \textbf{(D)}\ \frac{\sqrt{6}}{4}\qquad \textbf{(E)}\ \frac{\sqrt{6}-\sqrt{2}}{4}$


	\item (\textit{AHSME 1982}) Let $f(x)=|x-2|+|x-4|-|2x-6|$ for $2 \leq x\leq 8$. The sum of the largest and smallest values of $f(x)$ is 


 $\textbf {(A)}\ 1 \qquad  \textbf {(B)}\ 2 \qquad  \textbf {(C)}\ 4 \qquad  \textbf {(D)}\ 6 \qquad  \textbf {(E)}\ \text{none of these}$


	\item (\textit{AHSME 1982}) The number of pairs of positive integers $(x,y)$ which satisfy the equation $x^2+y^2=x^3$ is 


 $\textbf {(A)}\ 0 \qquad  \textbf {(B)}\ 1 \qquad  \textbf {(C)}\ 2 \qquad  \textbf {(D)}\ \text{not finite} \qquad  \textbf {(E)}\ \text{none of these}$


	\item (\textit{AHSME 1982}) In the adjoining figure, the triangle $ABC$ is a right triangle with $\angle BCA=90^\circ$. Median $CM$ is perpendicular to median $BN$, 
and side $BC=s$. The length of $BN$ is


 
                    \begin{figure}[H]
                        \centering
                        \includegraphics[width=0.5\textwidth]{Problem/AHSME_I_1982_21_Problem_Figure_1.png}
                    \end{figure}
                    
$\textbf{(A)}\ s\sqrt 2 \qquad  \textbf{(B)}\ \frac 32s\sqrt2 \qquad  \textbf{(C)}\ 2s\sqrt2 \qquad  \textbf{(D)}\ \frac{s\sqrt5}{2}\qquad \textbf{(E)}\ \frac{s\sqrt6}{2}$


	\item (\textit{AHSME 1982}) In a narrow alley of width $w$ a ladder of length $a$ is placed with its foot at point $P$ between the walls. 
Resting against one wall at $Q$, the distance $k$ above the ground makes a $45^\circ$ angle with the ground. 
Resting against the other wall at $R$, a distance $h$ above the ground, the ladder makes a $75^\circ$ angle with the ground. The width $w$ is equal to 


 $\textbf{(A)}\ a\qquad \textbf{(B)}\ RQ\qquad \textbf{(C)}\ k\qquad \textbf{(D)}\ \frac{h+k}{2}\qquad \textbf{(E)}\ h$


	\item (\textit{AHSME 1982}) The lengths of the sides of a triangle are consecutive integers, and the largest angle is twice the smallest angle. 
The cosine of the smallest angle is 


 $\textbf{(A)}\ \frac{3}{4}\qquad \textbf{(B)}\ \frac{7}{10}\qquad \textbf{(C)}\ \frac{2}{3}\qquad \textbf{(D)}\ \frac{9}{14}\qquad \textbf{(E)}\ \text{none of these}$


	\item (\textit{AHSME 1982}) In the adjoining figure, the circle meets the sides of an equilateral triangle at six points. If $AG=2, GF=13, FC=1$, 
and $HJ=7$, then $DE$ equals 


 
                    \begin{figure}[H]
                        \centering
                        \includegraphics[width=0.5\textwidth]{Problem/AHSME_I_1982_24_Problem_Figure_1.png}
                    \end{figure}
                    
$\textbf {(A)}\ 2\sqrt{22} \qquad  \textbf {(B)}\ 7\sqrt{3} \qquad  \textbf {(C)}\ 9 \qquad  \textbf {(D)}\ 10 \qquad  \textbf {(E)}\ 13$


	\item (\textit{AHSME 1982}) The adjacent map is part of a city: the small rectangles are blocks, and the paths in between are streets. 
Each morning, a student walks from intersection $A$ to intersection $B$, always walking along streets shown, 
and always going east or south. For variety, at each intersection where he has a choice, he chooses with 
probability $\frac{1}{2}$ whether to go east or south. Find the probability that through any given morning, he goes through $C$. 


 
                    \begin{figure}[H]
                        \centering
                        \includegraphics[width=0.5\textwidth]{Problem/AHSME_I_1982_25_Problem_Figure_1.png}
                    \end{figure}
                    
$\textbf{(A)}\ \frac{11}{32}\qquad \textbf{(B)}\ \frac{1}{2}\qquad \textbf{(C)}\ \frac{4}{7}\qquad \textbf{(D)}\ \frac{21}{32}\qquad \textbf{(E)}\ \frac{3}{4}$


	\item (\textit{AHSME 1982}) If the base $8$ representation of a perfect square is $ab3c$, where $a\ne 0$, then $c$ equals 


 $\textbf{(A)}\ 0\qquad  \textbf{(B)}\ 1 \qquad  \textbf{(C)}\ 3\qquad  \textbf{(D)}\ 4\qquad  \textbf{(E)}\ \text{not uniquely determined}$


	\item (\textit{AHSME 1982}) Suppose $z=a+bi$ is a solution of the polynomial equation 
 \[c_4z^4+ic_3z^3+c_2z^2+ic_1z+c_0=0,\] where $c_0, c_1, c_2, c_3, a,$ and $b$ are real constants and $i^2=-1.$ Which of the following must also be a solution? 


 $\textbf{(A)}\ -a-bi\qquad  \textbf{(B)}\ a-bi\qquad  \textbf{(C)}\ -a+bi\qquad  \textbf{(D)}\ b+ai \qquad  \textbf{(E)}\ \text{none of these}$


	\item (\textit{AHSME 1982}) A set of consecutive positive integers beginning with $1$ is written on a blackboard. One number is erased. The average (arithmetic mean) of the remaining numbers is $35\frac{7}{17}$. What number was erased? 


 $\textbf{(A)}\ 6\qquad  \textbf{(B)}\ 7 \qquad  \textbf{(C)}\ 8 \qquad  \textbf{(D)}\  9\qquad  \textbf{(E)}\ \text{cannot be determined}$


	\item (\textit{AHSME 1982}) Let $x,y,$ and $z$ be three positive real numbers whose sum is $1.$ If no one of these numbers is more than twice any other, 
then the minimum possible value of the product $xyz$ is 


 $\textbf{(A)}\ \frac{1}{32}\qquad \textbf{(B)}\ \frac{1}{36}\qquad \textbf{(C)}\ \frac{4}{125}\qquad \textbf{(D)}\ \frac{1}{127}\qquad \textbf{(E)}\ \text{none of these}$


	\item (\textit{AHSME 1982}) Find the units digit of the decimal expansion of 
 \[\left(15 + \sqrt{220}\right)^{19} + \left(15 + \sqrt{220}\right)^{82}.\]
$\textbf{(A)}\ 0\qquad  \textbf{(B)}\ 2\qquad  \textbf{(C)}\ 5\qquad  \textbf{(D)}\ 9\qquad  \textbf{(E)}\ \text{none of these}$


\end{enumerate}